% GENERATED FILE. Edit canonical lessons, then run npm run generate:books.
% canonical-source-hash: fnv1a64:85a761f55a45b38b
% canonical-lessons: HI-C539-tirth, HI-C539-dua, HI-C539-prarthna, HI-C539-kamna, HI-C539-thanna

\chapter{Wishes and Prayers}
\label{ch:hi-wishes-and-prayers}
\hlchaptermodality{539}

\begin{chapteropening}
\textbf{By the end of this chapter:} \emph{I can talk about what I wish and pray for and what I have resolved.}

Complete the last lesson of chapter 539: I can talk about what I wish and pray for and what I have resolved.
\end{chapteropening}

\section[\texorpdfstring{\hi{तीर्थ} (tīrth)}{तीर्थ (tīrth)}]{\hi{तीर्थ} (tīrth) — a place of pilgrimage}

\label{lesson:HI-C539-tirth}



\begin{quote}
Before the new one: say the Hindi for to be decorated, then the Hindi for to pick out, to pick over (grain).
\end{quote}

\subsection*{You'll want to know: \hi{तीर्थ}}
\textbf{\hi{तीर्थ}} — \emph{tīrth} — \textquotedblleft{}a place of pilgrimage\textquotedblright{}.

\begin{quote}
\hi{दादी} \hi{हर} \hi{साल} \hi{तीर्थ} \hi{पर} \hi{जाती} \hi{हैं।}
\end{quote}

\begin{quote}
\emph{Dādī har sāl tīrth par jātī haiṁ.}
\end{quote}

\begin{quote}
Every year my grandmother goes on a pilgrimage.
\end{quote}

Masculine.

\begin{cousinweb}
\textbf{\hi{तीर्थ}} (\emph{tīrth}) is a tatsama, taken whole from Sanskrit \emph{tīrtha}, \textquotedblleft{}a ford; a place of pilgrimage\textquotedblright{} (Monier-Williams): a holy place where one crosses over.
\end{cousinweb}

\begin{grammarlens}[title={What you've built}]
One more word for saying what happens if.
\end{grammarlens}

\subsection*{Your turn}
\begin{itemize}
\raggedright
  \item \emph{Say it:} \emph{tīrth}
  \item \emph{Say it:} \emph{tīrth}, once more
  \item \emph{Say it:} \emph{Dādī har sāl tīrth par jātī haiṁ.} — Every year my grandmother goes on a pilgrimage.
  \item \emph{Recall:} say the Hindi for to be decorated, then the Hindi for to pick out, to pick over (grain), then say \emph{tīrth} again
\end{itemize}

\begin{tcolorbox}[breakable,colback=teal!4,colframe=teal!35!black,title={Before you move on}]
What does \hi{तीर्थ} mean? (\textquotedblleft{}A place of pilgrimage\textquotedblright{}.) Say it once more.
\end{tcolorbox}
\section[\texorpdfstring{\hi{दुआ} (duā)}{दुआ (duā)}]{\hi{दुआ} (duā) — a prayer, a blessing (for someone)}

\label{lesson:HI-C539-dua}



\begin{quote}
Before the new one: say the Hindi for to pick out, to pick over (grain), then the Hindi for a place of pilgrimage.
\end{quote}

\subsection*{You'll want to know: \hi{दुआ}}
\textbf{\hi{दुआ}} — \emph{duā} — \textquotedblleft{}a prayer, a blessing (for someone)\textquotedblright{}.

\begin{quote}
\hi{दादी} \hi{ने} \hi{मेरी} \hi{सेहत} \hi{के} \hi{लिए} \hi{दुआ} \hi{की।}
\end{quote}

\begin{quote}
\emph{Dādī ne merī sehat ke liye duā kī.}
\end{quote}

\begin{quote}
Grandmother prayed for my health.
\end{quote}

Feminine.

\begin{cousinweb}
\textbf{\hi{दुआ}} (\emph{duā}) is usually traced to Arabic \emph{duʿā}, \textquotedblleft{}a call, a prayer\textquotedblright{}, through Persian.
\end{cousinweb}

\begin{grammarlens}[title={What you've built}]
One more word for saying what happens if.
\end{grammarlens}

\subsection*{Your turn}
\begin{itemize}
\raggedright
  \item \emph{Say it:} \emph{duā}
  \item \emph{Say it:} \emph{duā}, once more
  \item \emph{Say it:} \emph{Dādī ne merī sehat ke liye duā kī.} — Grandmother prayed for my health.
  \item \emph{Recall:} say the Hindi for to pick out, to pick over (grain), then the Hindi for a place of pilgrimage, then say \emph{duā} again
\end{itemize}

\begin{tcolorbox}[breakable,colback=teal!4,colframe=teal!35!black,title={Before you move on}]
What does \hi{दुआ} mean? (\textquotedblleft{}A prayer, a blessing (for someone)\textquotedblright{}.) Say it once more.
\end{tcolorbox}
\section[\texorpdfstring{\hi{प्रार्थना} (prārthnā)}{प्रार्थना (prārthnā)}]{\hi{प्रार्थना} (prārthnā) — a prayer; a request}

\label{lesson:HI-C539-prarthna}



\begin{quote}
Before the new one: say the Hindi for a place of pilgrimage, then the Hindi for a prayer, a blessing (for someone).
\end{quote}

\subsection*{You'll want to know: \hi{प्रार्थना}}
\textbf{\hi{प्रार्थना}} — \emph{prārthnā} — \textquotedblleft{}a prayer; a request\textquotedblright{}.

\begin{quote}
\hi{स्कूल} \hi{में} \hi{हर} \hi{सुबह} \hi{प्रार्थना} \hi{होती} \hi{है।}
\end{quote}

\begin{quote}
\emph{Skūl meṁ har subah prārthnā hotī hai.}
\end{quote}

\begin{quote}
There is a prayer at school every morning.
\end{quote}

Feminine.

\begin{cousinweb}
\textbf{\hi{प्रार्थना}} (\emph{prārthnā}) is a tatsama, taken whole from Sanskrit \emph{prārthanā}, \textquotedblleft{}a wish, a request, a prayer\textquotedblright{} (Monier-Williams).
\end{cousinweb}

\begin{grammarlens}[title={What you've built}]
One more word for saying what happens if.
\end{grammarlens}

\subsection*{Your turn}
\begin{itemize}
\raggedright
  \item \emph{Say it:} \emph{prārthnā}
  \item \emph{Say it:} \emph{prārthnā}, once more
  \item \emph{Say it:} \emph{Skūl meṁ har subah prārthnā hotī hai.} — There is a prayer at school every morning.
  \item \emph{Recall:} say the Hindi for a place of pilgrimage, then the Hindi for a prayer, a blessing (for someone), then say \emph{prārthnā} again
\end{itemize}

\begin{tcolorbox}[breakable,colback=teal!4,colframe=teal!35!black,title={Before you move on}]
What does \hi{प्रार्थना} mean? (\textquotedblleft{}A prayer; a request\textquotedblright{}.) Say it once more.
\end{tcolorbox}
\section[\texorpdfstring{\hi{कामना} (kāmnā)}{कामना (kāmnā)}]{\hi{कामना} (kāmnā) — a wish (for someone's good)}

\label{lesson:HI-C539-kamna}



\begin{quote}
Before the new one: say the Hindi for a prayer, a blessing (for someone), then the Hindi for a prayer.
\end{quote}

\subsection*{You'll want to know: \hi{कामना}}
\textbf{\hi{कामना}} — \emph{kāmnā} — \textquotedblleft{}a wish (for someone's good)\textquotedblright{}.

\begin{quote}
\hi{हम} \hi{आपकी} \hi{अच्छी} \hi{सेहत} \hi{की} \hi{कामना} \hi{करते} \hi{हैं।}
\end{quote}

\begin{quote}
\emph{Ham āpkī acchī sehat kī kāmnā karte haiṁ.}
\end{quote}

\begin{quote}
We wish you good health.
\end{quote}

Feminine.

\begin{cousinweb}
\textbf{\hi{कामना}} (\emph{kāmnā}) is a tatsama, taken whole from Sanskrit \emph{kāmanā}, \textquotedblleft{}wish, desire\textquotedblright{} (Monier-Williams).
\end{cousinweb}

\begin{grammarlens}[title={What you've built}]
One more word for saying what happens if.
\end{grammarlens}

\subsection*{Your turn}
\begin{itemize}
\raggedright
  \item \emph{Say it:} \emph{kāmnā}
  \item \emph{Say it:} \emph{kāmnā}, once more
  \item \emph{Say it:} \emph{Ham āpkī acchī sehat kī kāmnā karte haiṁ.} — We wish you good health.
  \item \emph{Recall:} say the Hindi for a prayer, a blessing (for someone), then the Hindi for a prayer, then say \emph{kāmnā} again
\end{itemize}

\begin{tcolorbox}[breakable,colback=teal!4,colframe=teal!35!black,title={Before you move on}]
What does \hi{कामना} mean? (\textquotedblleft{}A wish (for someone's good)\textquotedblright{}.) Say it once more.
\end{tcolorbox}
\section[\texorpdfstring{\hi{ठानना} (ṭhānnā)}{ठानना (ṭhānnā)}]{\hi{ठानना} (ṭhānnā) — to resolve; to make up one's mind}

\label{lesson:HI-C539-thanna}



\begin{quote}
Before the new one: say the Hindi for a prayer, then the Hindi for a wish (for someone's good).
\end{quote}

\subsection*{You'll want to know: \hi{ठानना}}
\textbf{\hi{ठानना}} — \emph{ṭhānnā} — \textquotedblleft{}to resolve; to make up one's mind\textquotedblright{}.

\begin{quote}
\hi{मीरा} \hi{ने} \hi{ठान} \hi{लिया} \hi{कि} \hi{वह} \hi{तैरना} \hi{सीखेगी।}
\end{quote}

\begin{quote}
\emph{Mīrā ne ṭhān liyā ki vah tairnā sīkhegī.}
\end{quote}

\begin{quote}
Mira made up her mind that she would learn to swim.
\end{quote}

Past, with \hi{ने} (he or she …): \hi{उसने} \hi{ठान} \hi{लिया}, \emph{usne ṭhān liyā}.

\begin{cousinweb}
\textbf{\hi{ठानना}} (\emph{ṭhānnā}) is a tadbhava, inherited from Sanskrit \emph{sthāna}, \textquotedblleft{}a firm stance, a place\textquotedblright{} (CDIAL 13753). CDIAL gives the Hindi verb as \textquotedblleft{}to settle, determine\textquotedblright{}.
\end{cousinweb}

\begin{grammarlens}[title={What you've built}]
One more word for saying what happens if.
\end{grammarlens}

\subsection*{Your turn}
\begin{itemize}
\raggedright
  \item \emph{Say it:} \emph{ṭhānnā}
  \item \emph{Say it:} \emph{ṭhānnā}, once more
  \item \emph{Say it:} \emph{Mīrā ne ṭhān liyā ki vah tairnā sīkhegī.} — Mira made up her mind that she would learn to swim.
  \item \emph{Recall:} say the Hindi for a prayer, then the Hindi for a wish (for someone's good), then say \emph{ṭhānnā} again
\end{itemize}

\begin{tcolorbox}[breakable,colback=teal!4,colframe=teal!35!black,title={Before you move on}]
What does \hi{ठानना} mean? (\textquotedblleft{}To resolve; to make up one's mind\textquotedblright{}.) Say it once more.
\end{tcolorbox}
